\chapter{Conclusions and Future Work}
\label{chap:conc-trab-futuro}

\section{Main Conclusions}
\label{sec:conc-princ}

The work described in this report addresses a concrete problem of
the \ac{SEGAL} laboratory: how to refactor a large, several-years-old
Jakarta EE \ac{API} without silently breaking the \ac{EPOS} clients
that depend on it. The answer is a three-stage \ac{CI}/\ac{CD}-driven
testing framework that, for every commit, compares the new build
against two baselines (the immediate parent commit and a fixed
long-term reference) both at the level of internal Java classes and
at the level of \ac{HTTP} responses.

A few conclusions stand out. \emph{Differential testing fits a
refactor naturally}: asserting that the new build produces the same
output as a known-good one sidesteps the brittle fixture problem and
made the framework useful from the very first commit it ran on.
\emph{Two layers of comparison are better than one}: the code-level
harness catches behavioural changes in focused helper classes with
fast feedback, the \ac{HTTP}-level harness catches integration
regressions visible only against a live database, and neither one
alone covers the spectrum the laboratory cares about. \emph{A
self-hosted pipeline is not optional here}: only a runner inside the
laboratory's network can exercise the \ac{API} end-to-end against
the production database, so self-managing both the GitLab control
plane and the GitLab Runner host followed from the data-locality
constraint, not from a preference. \emph{The two prototypes paid for
themselves}: each answered a small, well-bounded set of questions in
an environment where mistakes were cheap, sparing the production
code base from them. \emph{The framework reveals what a human review
would miss}: the most representative example, discussed in
Section~\ref{sec:reports}, is the named-constants refactor of the
latitude/longitude bounds that no amount of code review would
plausibly have caught at the laboratory's commit rate.

\section{Future Work}
\label{sec:trab-futuro}

Several extensions to the framework were deliberately left out of
scope. The most immediate are the introduction of \textbf{BRIN}
indexes on the date columns of \texttt{rinex\_file} (and the
companion \emph{performance-regression subsystem} that would record
wall-clock durations against the \texttt{F} fixture to justify them
empirically), the \textbf{PostGIS migration} of the spatial helpers
in the \texttt{Geometry} package (with the same differential
framework used to verify byte-for-byte equivalence), and the
\textbf{automatic generation of \ac{HTTP} test cases} from the
\ac{OpenAPI} document shipped inside the \ac{WAR} (so newly added
endpoints are exercised by the diff from the moment they are
merged). The framework would also benefit from a \textbf{schema diff}
of the \ac{OpenAPI} document itself across builds (to surface
contract regressions distinct from behavioural ones), from a
\textbf{snapshot test} layer for endpoints that should remain stable
across releases (such as \texttt{/api/ping}), and from
\textbf{containerization of the production deployment} to remove
the residual gap between the \ac{WAR} that was tested and the
\ac{WAR} that ships. Finally, the harness code itself could be
\textbf{generalised into a small library} reusable by other Jakarta
EE services of the laboratory, with the application-specific bits
factored into a clean configuration interface; orthogonal to that,
several \textbf{operational aspects} (transient-error retries, clean
termination of orphan Payara Micro processes, automated provisioning
of the runner host, a long-term dashboard on top of the diff
reports) deserve more attention than the present project could
afford.

\bigskip

Overall, the framework delivered by this project addresses the
problem that motivated it, has already proved its value on a real
regression caught during the refactor, and lays the groundwork for
the laboratory to continue evolving the \ac{API} with a much higher
degree of confidence than was previously possible. The items listed
in this section are natural next steps; none of them invalidates
the design that was put in place.
